\begin{explanationbox}
\textbf{\textcolor{CExplanation}{一句话直觉}}

等比数列求和，本质是利用“乘以公比后错位相减”消去中间项，化繁为简。

\medskip
\textbf{\textcolor{CExplanation}{核心拆解}}
\begin{description}[style=nextline, leftmargin=2.8em, labelsep=0.5em, itemsep=0.45em]
    \item[\textbf{要点一（公比作用）：}] 公比 $q$ 控制数列的放大、缩小或正负交替。当 $|q|>1$ 时和增长极快，$|q|<1$ 时和趋近常数。
    \item[\textbf{要点二（错位消去）：}] 将 $S_n$ 与 $qS_n$ 对齐相减，绝大多数项成对抵消，只留下首尾两项。
    \item[\textbf{要点三（分类讨论）：}] $q=1$ 时数列退化为常数列，求和就是 $n$ 个 $a_1$ 相加，直接写 $na_1$。
\end{description}

\medskip
\textbf{\textcolor{CExplanation}{几何本质（相似缩放）}}

等比数列的各项可以看作相似图形边长依次缩放 $q$ 倍，前 $n$ 项和对应缩放过程的累积总量。

\medskip
\textbf{\textcolor{CExplanation}{代数意义（方程转换）}}

求和公式通过构造一次方程 $S_n - q S_n = a_1 - a_n q$ 解出，体现了递推关系转化为封闭表达式的思想。

\medskip
\textbf{\textcolor{CExplanation}{考点价值（题型归纳）}}
\begin{itemize}[leftmargin=*, itemsep=0.5em, topsep=0.4em]
    \item \textbf{考法一（直接求值）：} 已知 $a_1,q,n$ 中的三个量，直接用公式计算 $S_n$。
    \item \textbf{考法二（反求参数）：} 已知 $S_n$ 及部分条件，反解 $q$ 或 $n$，注意可能有多解。
    \item \textbf{考法三（综合递推）：} 与错位相减法结合，求解数列通项或用于放缩。
\end{itemize}

\medskip
\textbf{\textcolor{CExplanation}{顿悟点}}

\textbf{“错位相减”不是技巧，而是等比数列定义（后项比前项恒为 $q$）的直接代数表达。}

\medskip
\textbf{\textcolor{CExplanation}{使用场景}}
\begin{itemize}[leftmargin=*, itemsep=0.5em, topsep=0.4em]
    \item \textbf{场景一（复利模型）：} 本金加利息按固定利率滚动，多次存款的终值和就是等比求和。
    \item \textbf{场景二（无穷级数）：} 无穷递减等比级数（如 $1+\frac12+\frac14+\dots$）直接套 $S=\frac{a_1}{1-q}$。
\end{itemize}
\end{explanationbox}
